\section{Introduction}
A defining and growing cost of operating a full node in Bitcoin~\cite{nakamoto2008bitcoin} or
Ethereum is storing the complete transaction history. This cost is \emph{unbounded}, it
only ever increases, and it is paid by \emph{every} full node, because the classical
security model replicates the whole chain everywhere. The practical consequence is
centralization: as history accumulates, fewer participants can afford to validate
independently, and the network drifts toward a small set of storage-rich operators. Energy
is the second well-known cost: PoW chains expend electricity solely to decide who appends the
next block, work that is intrinsically discarded.

The magnitude is not hypothetical. Production chains accumulate history into the hundreds of
gigabytes and beyond, and the requirement that every validating node retain all of it is widely
recognized as a scalability and decentralization bottleneck, motivating an entire line of
coded-storage work~\cite{yang2024scaling,perard2018erasure}. Yet most such work treats coding as
an optimization bolted onto an unchanged consensus. We instead ask what a chain looks like when
fragmentation is the \emph{default} and when block-production rights are derived from the very act
of storing fragments.

We take the position that the scarce, security-relevant resource of a blockchain should be
\emph{storage of the network's own data}, not wasted hashing, and that \emph{no single node
should need the entire history} to participate. \textbf{ShadowLedger} is our realization of
this position. Its thesis has two parts. First, block \emph{history} is erasure-coded and
scattered across the network so each node retains only a fraction, yet any sufficiently large
set of nodes can reconstruct any block on demand. Second, the right to produce blocks is
earned by \emph{proving storage} of assigned fragments over a bonded validator set, rather
than by burning energy; the only PoW present is a one-shot Sybil gate and a small onboarding
faucet, never the block-ordering mechanism.

Realizing this is not free of tension. If history is not stored whole, validators must still
be able to recover it, or a malicious producer could withhold data; this is the
\emph{data-availability} problem~\cite{albassam2021fraud,yu2020coded}. If block production is
permissionless, the system must resist Sybil flooding and equivocation. And if storage
replaces hashing, the protocol must measure and reward it without a centralized auditor. We
address each of these and report what a working implementation achieves and what it does not.

\paragraph{Contributions.}
\begin{enumerate}
  \item An architecture that fragments \emph{history} (block bodies) with Reed--Solomon coding
  and rendezvous placement while keeping \emph{state} (balances, nonces, contract storage)
  materialized locally, the \emph{state-versus-history} separation that reconciles ``no node
  stores everything'' with fast stateless validation (\S\ref{sec:history}).
  \item A Proof-of-Storage consensus with a bonded on-chain validator registry, a one-shot
  registration PoW as a Sybil gate, HRW leader election with a round-based liveness fallback,
  equivocation slashing, a heaviest-chain reorg engine, and depth-based finality
  (\S\ref{sec:consensus}, \S\ref{sec:fork}). Leader election is weighted by
  on-chain proven storage (validators submit storage-proof records that drive election odds); what
  remains specified-not-built is \emph{bond} slashing for missed proofs (the prototype uses weight
  decay) and storage-weighted \emph{fork choice} (\S\ref{sec:audit}, \S\ref{sec:disc}).
  \item A capital-free onboarding mechanism (a PoW faucet) and a deterministic contract VM with
  a Solidity-like language, both anchored to committed block data and thus verifiable without
  holding the fragmented body (\S\ref{sec:econ}, \S\ref{sec:vm}).
  \item An analytical storage model ($O(1/n)$ per-node history) and an empirical evaluation of a
  live Go prototype, with an explicit threat model and an honest account of open problems
  (\S\ref{sec:threat}--\S\ref{sec:disc}).
\end{enumerate}

The paper is organized as follows. \S\ref{sec:bg} and \S\ref{sec:related} review the primitives
(erasure coding, rendezvous hashing, data availability, storage-based consensus) and position
ShadowLedger against prior coded blockchains, sharding, and proof-of-storage systems.
\S\ref{sec:overview} states the design goals and gives the node architecture. \S\ref{sec:history}
develops the erasure-coded history and its storage model; \S\ref{sec:consensus}--\ref{sec:fork}
the Proof-of-Storage consensus, storage auditing, fork choice, and finality; \S\ref{sec:econ} the
economics and capital-free onboarding; and \S\ref{sec:vm} the contract VM. \S\ref{sec:threat}
presents the threat model, \S\ref{sec:impl} the implementation, and \S\ref{sec:eval} the
evaluation. \S\ref{sec:disc}--\ref{sec:rationale} discuss limitations, applicability, and design
rationale before \S\ref{sec:concl} concludes.
